% Bagian: sets-functions-relations
% Bab: arithmetization
% Subbagian: Dari N ke Z

\documentclass[../../../include/open-logic-section]{subfiles}


\begin{document}

\olfileid[id]{sfr}{arith}{int}
\olsection{Dari $\Nat$ ke $\Int$}

Berikut dua pengamatan dasar:
	\begin{enumerate}
		\item Setiap bilangan bulat dapat ditulis dalam bentuk $n - m$, dengan $n, m \in\Nat$.
		\item Informasi yang dikodekan dalam suatu ekspresi $n - m$ dapat pula dikodekan oleh pasangan terurut $\tuple{n, m}$.
	\end{enumerate}
Kita sudah mengetahui bahwa pasangan-pasangan terurut bilangan asli merupakan !!{element}s dari $\Nat^2$. Kita juga mengandaikan bahwa kita memahami $\Nat$. Maka, berdasarkan dua pengamatan tadi, berikut sebuah usulan naif: \emph{mari kita perlakukan bilangan bulat sebagai pasangan terurut bilangan asli}.

Sebenarnya, usulan ini terlalu naif. Kita tentu menginginkan $0- 2 = 4 - 6$. Akan tetapi, jelas bahwa $\tuple{0, 2 }\neq \tuple{4, 6}$. Jadi, kita tidak dapat begitu saja mengatakan bahwa $\Nat^2$ adalah himpunan bilangan bulat.

Dengan menggeneralisasi masalah tadi, yang kita inginkan adalah:
	$$a - b = c - d \text{ jika dan hanya jika }a + d = c + b$$
(Sudah semestinya jelas bahwa begitulah bilangan bulat \emph{dimaksudkan} berperilaku: cukup tambahkan $b$ dan $d$ pada kedua ruas.) Cara mudah untuk menjamin perilaku ini ialah dengan mendefinisikan suatu relasi ekuivalensi antara pasangan-pasangan terurut, $\Intequiv$, sebagai berikut:
$$\tuple{ a, b } \Intequiv \tuple{c, d}\text{ jika dan hanya jika }a + d = c + b$$
Sekarang kita harus menunjukkan bahwa relasi ini merupakan relasi ekuivalensi.
\begin{prop} $\Intequiv$ merupakan relasi ekuivalensi.
\end{prop}
	\begin{proof}
	Kita harus menunjukkan bahwa $\Intequiv$ bersifat refleksif, simetris, dan transitif.

	\emph{Refleksivitas:} Jelas bahwa $\tuple{a, b} \Intequiv \tuple{a, b}$, karena $a + b = b + a$.

	\emph{Simetri:} Andaikan $\tuple{a, b} \Intequiv \tuple{c, d}$, sehingga $a + d = c + b$. Maka $c + b = a + d$, sehingga $\tuple{c, d} \Intequiv \tuple{a, b}$.

	\emph{Transitivitas:} Andaikan $\tuple{a, b} \Intequiv \tuple{c, d}\Intequiv \tuple{m, n}$. Maka $a + d = c + b$ dan $c + n = m + d$. Jadi $a + d + c + n = c + b + m + d$, dan dengan demikian $a + n = m + b$. Oleh karena itu, $\tuple{a, b } \Intequiv \tuple{m, n}$.
\end{proof}

Sekarang kita dapat menggunakan relasi ekuivalensi ini untuk membentuk kelas-kelas ekuivalensi:

\begin{defn}
		Bilangan bulat adalah kelas-kelas ekuivalensi, di bawah $\Intequiv$, dari pasangan-pasangan terurut bilangan asli; yakni, $\Int = \equivclass{\Nat^2}{\Intequiv}$.
\end{defn}

Orang mungkin mempunyai berbagai reaksi \emph{filosofis} terhadap definisi stipulatif ini. Namun, sebelum mempertimbangkan reaksi-reaksi tersebut, ada baiknya kita melanjutkan beberapa rincian teknis.

Setelah menyatakan apa itu bilangan bulat, kita perlu mendefinisikan fungsi dan relasi dasar padanya. Mari kita tulis $\equivrep{m, n}{\Intequiv}$ untuk kelas ekuivalensi di bawah $\Intequiv$ yang memuat $\tuple{m, n}$ sebagai !!a{element}.\footnote{Catatan: dengan menggunakan notasi yang diperkenalkan dalam \olref[sfr][rel][eqv]{def:equivalenceclass}, kita akan menulis $\equivrep{\tuple{m,n}}{\Intequiv}$ untuk hal yang sama. Namun, bentuk itu sedikit lebih sulit dibaca.} Yakni:
	$$\equivrep{m, n}{\Intequiv} = \Setabs{\tuple{a, b} \in \Nat^2}{\tuple{a, b}\Intequiv \tuple{m, n}}$$
Sekarang kita memberikan beberapa definisi:
	\begin{align*}
		\equivrep{a, b}{\Intequiv} + \equivrep{c, d}{\Intequiv} &= \equivrep{a + c, b + d}{\Intequiv}\\
		\equivrep{a, b}{\Intequiv} \times \equivrep{c, d}{\Intequiv} &= \equivrep{a c + b  d, a  d + b c}{\Intequiv}\\
		\equivrep{a, b}{\Intequiv} \leq \equivrep{c, d}{\Intequiv} &\text{ jika dan hanya jika }a + d \leq b + c
	\end{align*}
(Sebagaimana lazimnya, saya menggunakan `$ab$' untuk menyatakan `$(a \times b)$', semata-mata agar aksioma-aksioma tersebut lebih mudah dibaca.) Sekarang kita perlu memastikan bahwa definisi-definisi ini berperilaku sebagaimana \emph{semestinya}. Menjabarkan apa artinya dan memeriksanya satu per satu cukup melelahkan; rinciannya kita tempatkan dalam \olref[check]{sec}. Namun, intinya: semuanya berhasil!

Masih ada satu hal terakhir. Kita telah membangun bilangan bulat dengan
menggunakan bilangan asli. Akan tetapi, ini berarti bahwa bilangan asli
\emph{bukan bilangan bulat itu sendiri}. Kita akan kembali membahas arti
filosofisnya dalam \olref[ref]{sec}. Namun, dari sudut pandang teknis murni,
kita memerlukan suatu cara agar dapat memperlakukan bilangan asli
\emph{sebagai} bilangan bulat. Gagasannya cukup mudah: untuk setiap $n \in
\Nat$, kita menetapkan bahwa $n_\Int = \equivrep{n, 0}{\Intequiv}$. Kita
perlu memastikan bahwa definisi ini berperilaku baik, yaitu bahwa untuk
sembarang $m, n \in \Nat$
\begin{align*}
	(m + n)_\Int &= m_\Int + n_\Int\\
	(m \times n)_\Int &= m_\Int \times n_\Int\\
	m \leq n &\liff m_\Int \leq n_\Int
\end{align*}
Namun, semua ini cukup langsung. Sebagai contoh, untuk menunjukkan bahwa
kondisi kedua berlaku, kita cukup memanfaatkan perilaku bilangan asli dan
bernalar sebagai berikut:
%\begin{proof}
%	Dengan memanfaatkan perilaku bilangan asli, jelas bahwa:
	\begin{align*}
%		(m + n)_\Int  = \equivrep{m + n, 0}{\Intequiv} = \equivrep{m + n, 0 + 0}{\Intequiv} = \equivrep{m, 0}{\Intequiv} + \equivrep{n, 0}{\Intequiv} = m_\Int + n_\Int\\
		(m \times n)_\Int  &= \equivrep{m \times n, 0}{\Intequiv} \\
		&= \equivrep{m \times n + 0 \times 0, m \times 0 + 0 \times n}{\Intequiv} \\
		&= \equivrep{m, 0}{\Intequiv} \times \equivrep{n, 0}{\Intequiv} \\
		&= m_\Int \times n_\Int
	\end{align*}
%\end{proof}
Sebagai latihan, silakan pastikan bahwa kedua kondisi lainnya berlaku.
\begin{prob}
	Tunjukkan bahwa $(m + n)_\Int = m_\Int + n_\Int$ dan $m \leq n \liff m_\Int \leq n_\Int$, untuk sembarang $m, n \in \Nat$.
\end{prob}
\end{document}

